\section{Семинар 6. Комбинаторные задачи, отображения и беспорядки}
\noeqref{eq:div_tower_total, eq:grid_paths, eq:tri_quad_sum, eq:map_all}
\noeqref{eq:map_strict_inc, eq:map_weak_inc, eq:disjoint_subsets_binom}
\noeqref{eq:diag_sub_method, eq:diag_direct_method, eq:diag_closed_form}
\noeqref{eq:derange_pie, eq:derange_rec_1, eq:derange_rec_2}

\begin{remark}[Обзор темы]
    Шестой семинар посвящён решению задач перечислительной комбинаторики 
    различных типов: применению правил суммы и произведения в теоретико-числовых 
    и геометрических сюжетах, подсчёту числа произвольных и монотонных 
    отображений между конечными множествами, перечислению пар 
    непересекающихся подмножеств и непересекающихся диагоналей выпуклого 
    многоугольника, а также фундаментальной задаче о беспорядках, решаемой 
    методом включений-исключений и методом рекуррентных соотношений.
\end{remark}

\subsection{Комбинаторный анализ в задачах делимости и геометрии}

\begin{problem}[\textbf{Делимость степенной башни}]
    Дано множество $A = \{1, 2, 3, 4\}$. Из множества $A$ выбираются с 
    возможными повторениями три числа $a, b, c \in A$. Сколько существует 
    упорядоченных троек $(a, b, c) \in A^3$, для которых значение 
    степенной башни $a^{b^c} = a^{(b^c)}$ делится нацело на $4$?
\end{problem}

\begin{proof}[Решение]~\\
    \begin{quote}
        Всего существует $\abs{A}^3 = 4^3 = 64$ упорядоченные тройки 
        $(a, b, c)$. Поскольку показатели $b, c \in \{1, 2, 3, 4\}$ 
        являются натуральными числами, показатель степени $k = b^c \ge 1$ 
        также натурален. Следовательно, делимость степени $a^k$ на $4$ 
        зависит прежде всего от чётности основания $a \in \{1, 2, 3, 4\}$. 
        Разберём все возможные значения основания $a$:
        \begin{enumerate}
            \item \textbf{Случай $a \in \{1, 3\}$ (нечётное основание):} 
            Любая натуральная степень нечётного числа нечётна и не может 
            делиться на $4$. Этот случай даёт $0$ решений.
            \item \textbf{Случай $a = 4$:} 
            Так как $b^c \ge 1$ при любых $b, c \in A$, число 
            $4^{b^c} = 4 \cdot 4^{b^c - 1}$ всегда кратно $4$ независимо 
            от выбора $b$ и $c$. Поскольку $b \in A$ и $c \in A$ выбираются 
            произвольно по $4$ способа каждое, по правилу умножения получаем:
            \[
                N_1 = 1 \cdot 4 \cdot 4 = 4^2 = 16 \text{ вариантов.}
            \]
            \item \textbf{Случай $a = 2$:} 
            Степень двойки $2^{b^c}$ делится на $4 = 2^2$ тогда и только 
            тогда, когда показатель степени удовлетворяет неравенству 
            $b^c \ge 2$.
            \begin{itemize}
                \item Если $b = 1$, то при любом $c \in A$ имеем 
                $1^c = 1 < 2$, откуда $2^{1^c} = 2 \not\vdots 4$.
                \item Если же $b \ge 2$ (то есть $b \in \{2, 3, 4\}$ --- 
                $3$ варианта), то при любом $c \in \{1, 2, 3, 4\}$ ($4$ 
                варианта) выполняется $b^c \ge 2^1 = 2$, а значит, 
                $2^{b^c} \vdots 4$.
            \end{itemize}
            По правилу умножения этот случай даёт:
            \[
                N_2 = 1 \cdot 3 \cdot 4 = 12 \text{ вариантов.}
            \]
        \end{enumerate}
        Поскольку случаи $a = 4$ и $a = 2$ не пересекаются, по правилу 
        сложения находим общее количество подходящих троек:
        \begin{equation} \label{eq:div_tower_total}
            N = N_1 + N_2 = 16 + 12 = 28 \text{ вариантов.} \tag{1}
        \end{equation}
        \textbf{Ответ:} $28$ вариантов.
    \end{quote}
\end{proof}

\begin{problem}[\textbf{Кратчайшие пути на клетчатой плоскости}]
    Рассмотрим целочисленную решётку $\Z^2$ на плоскости. Найти количество 
    кратчайших путей по линиям сетки из начального узла $(0, 0)$ в узел 
    $(m, n)$, где $m, n \in \N \cup \{0\}$.
\end{problem}

\begin{proof}[Решение]~\\
    \begin{quote}
        По линиям клетчатой сетки из каждого узла $(x, y)$ можно сделать 
        шаг длины $1$ в четырёх направлениях: вправо ($\to$), влево ($\leftarrow$), 
        вверх ($\uparrow$) или вниз ($\downarrow$). Расстояние между 
        узлами $(0, 0)$ и $(m, n)$ в манхэттенской метрике равно $m + n$. 
        Чтобы путь был \emph{кратчайшим}, на каждом шаге мы должны 
        приближаться к целевой точке $(m, n)$, то есть двигаться только 
        \textbf{вправо} (увеличивая абсциссу на $1$) или \textbf{вверх} 
        (увеличивая ординату на $1$).

        \begin{figure}[H]
            \centering
            \begin{tikzpicture}[scale=0.85]
                % Сетка m x n (для примера m = 5, n = 3)
                \draw[step=1cm, gray!50, thin] (0, 0) grid (5, 3);

                % Оси и граничные подписи
                \fill (0, 0) circle (2.5pt) node[below left] {$(0, 0)$};
                \fill (5, 3) circle (2.5pt) node[above right] {$(m, n)$};

                % Пример кратчайшего пути: П, В, П, П, В, П, В, П
                \draw[->, blue!80!black, ultra thick] 
                    (0, 0) -- (1, 0) -- (1, 1) -- (3, 1) -- (3, 2) 
                    -- (4, 2) -- (4, 3) -- (5, 3);

                % Выделение узлов на пути
                \foreach \x/\y in {1/0, 1/1, 2/1, 3/1, 3/2, 4/2, 4/3} {
                    \fill[blue!80!black] (\x, \y) circle (1.8pt);
                }

                % Иконка направлений справа
                \begin{scope}[shift={(6.8, 1.5)}, scale=0.6]
                    \draw[<->, thick] (-1, 0) -- (1, 0);
                    \draw[<->, thick] (0, -1) -- (0, 1);
                    \node[right] at (1.1, 0) {допустимы только $\to$ и $\uparrow$};
                \end{scope}
            \end{tikzpicture}
            \caption{Кратчайший путь на решётке из $(0, 0)$ в $(m, n)$ 
            состоит ровно из $m$ шагов вправо и $n$ шагов вверх.}
        \end{figure}

        Таким образом, любой кратчайший путь имеет длину $m + n$ и 
        однозначно кодируется словом длины $m + n$ из двух символов 
        $\{\text{П}, \text{В}\}$, в котором символ «П» (шаг вправо) 
        встречается ровно $m$ раз, а символ «В» (шаг вверх) --- ровно $n$ 
        раз. Количество таких слов равно числу способов выбрать $n$ позиций 
        для шагов вверх (или $m$ позиций для шагов вправо) из $m + n$ 
        доступных шагов:
        \begin{equation} \label{eq:grid_paths}
            N = C_{m+n}^n = C_{m+n}^m = \frac{(m + n)!}{m! \, n!}. \tag{2}
        \end{equation}
        \textbf{Ответ:} $C_{m+n}^n$.
    \end{quote}
\end{proof}

\begin{problem}[\textbf{Треугольники и четырёхугольники на параллельных прямых}]
    На плоскости даны две параллельные прямые $l_1 \parallel l_2$. На 
    первой прямой отмечены $10$ различных точек, а на второй --- $12$ 
    различных точек. Найти количество невырожденных треугольников и 
    четырёхугольников с вершинами в отмеченных точках, а также их общее 
    число.
\end{problem}

\begin{proof}[Решение]~\\
    \begin{quote}
        Никакие три вершины невырожденного многоугольника (треугольника или 
        четырёхугольника) не могут лежать на одной прямой.

        \textbf{1. Подсчёт числа треугольников ($N_{\triangle}$):}\\
        Из трёх вершин треугольника две должны лежать на одной из 
        параллельных прямых, а третья вершина --- на другой прямой. 
        Получаем два непересекающихся случая:
        \begin{itemize}
            \item $2$ вершины на первой прямой ($C_{10}^2$ способов) и $1$ 
            вершина на второй прямой ($C_{12}^1 = 12$ способов): всего 
            $C_{10}^2 \cdot 12$ треугольников;
            \item $1$ вершина на первой прямой ($C_{10}^1 = 10$ способов) и 
            $2$ вершины на второй прямой ($C_{12}^2$ способов): всего 
            $C_{12}^2 \cdot 10$ треугольников.
        \end{itemize}
        Всего треугольников: $N_{\triangle} = C_{10}^2 \cdot 12 + 
        C_{12}^2 \cdot 10 = 45 \cdot 12 + 66 \cdot 10 = 540 + 660 = 1200$.

        \textbf{2. Подсчёт числа четырёхугольников ($N_{\square}$):}\\
        Чтобы никакие три вершины четырёхугольника не лежали на одной 
        прямой, необходимо выбрать ровно $2$ вершины на первой прямой 
        ($C_{10}^2$ способов) и ровно $2$ вершины на второй прямой 
        ($C_{12}^2$ способов). Любая такая четвёрка точек однозначно 
        образует выпуклый четырёхугольник (трапецию или параллелограмм). 
        По правилу умножения:
        \[
            N_{\square} = C_{10}^2 \cdot C_{12}^2 = 45 \cdot 66 = 2970.
        \]

        \textbf{3. Суммарное количество фигур:}\\
        По правилу сложения общее число треугольников и четырёхугольников 
        равно:
        \begin{equation} \label{eq:tri_quad_sum}
            N_{\text{общ}} = N_{\triangle} + N_{\square} = 
            C_{10}^2 \cdot 12 + C_{12}^2 \cdot 10 + C_{10}^2 \cdot C_{12}^2 = 
            1200 + 2970 = 4170. \tag{3}
        \end{equation}
    \end{quote}
\end{proof}

\subsection{Подсчёт числа отображений и монотонных функций}

\begin{problem}[\textbf{Произвольные и монотонные отображения}]
    Пусть даны два конечных упорядоченных множества 
    $X = \{1, 2, \dots, m\}$ и $Y = \{1, 2, \dots, n\}$. Найти:
    \begin{enumerate}
        \item общее число всех отображений $f: X \to Y$;
        \item число \textbf{строго возрастающих} отображений $f: X \to Y$ 
        (при условии $m \le n$);
        \item число \textbf{нестрого возрастающих} (неубывающих) 
        отображений $f: X \to Y$.
    \end{enumerate}
\end{problem}

\begin{proof}[Решение]~\\
    \begin{quote}
        Каждое отображение $f: \{1, \dots, m\} \to \{1, \dots, n\}$ 
        взаимно однозначно задаётся упорядоченным кортежем своих значений 
        $\bigl(f(1), f(2), \dots, f(m)\bigr) \in Y^m$.

        \textbf{1. Все отображения $f: \{1, \dots, m\} \to \{1, \dots, n\}$:}\\
        Для каждого из $m$ элементов области определения $i \in \{1, \dots, m\}$ 
        образ $f(i) \in \{1, \dots, n\}$ выбирается независимо любым из $n$ 
        способов. Следовательно, число всех отображений равно числу 
        размещений с повторениями из $n$ элементов по $m$:
        \begin{equation} \label{eq:map_all}
            N_{\text{все}} = \bar{A}_n^m = n^m. \tag{4}
        \end{equation}

        \textbf{2. Строго возрастающие отображения ($1 \le i < j \le m \implies 
        f(i) < f(j)$):}\\
        Такое отображение задаётся цепочкой строгих неравенств:
        \[
            1 \le f(1) < f(2) < \dots < f(m) \le n.
        \]
        Поскольку все $m$ значений $f(1), \dots, f(m)$ обязаны быть попарно 
        различными, при $m > n$ таких отображений не существует ($0$). При 
        $m \le n$ любое $m$-элементное подмножество $\{y_1, y_2, \dots, y_m\} 
        \subseteq \{1, \dots, n\}$ можно упорядочить по возрастанию 
        \emph{единственным} образом: $y_1 < y_2 < \dots < y_m$, положив 
        $f(i) = y_i$. Тем самым установлена биекция между строго 
        возрастающими отображениями и $m$-элементными подмножествами 
        $n$-множества (сочетаниями без повторений):
        \begin{equation} \label{eq:map_strict_inc}
            N_{\text{стр}} = C_n^m = \binom{n}{m}. \tag{5}
        \end{equation}

        \textbf{3. Нестрого возрастающие отображения ($1 \le i < j \le m \implies 
        f(i) \le f(j)$):}\\
        Такое отображение определяется цепочкой нестрогих неравенств:
        \[
            1 \le f(1) \le f(2) \le \dots \le f(m) \le n.
        \]
        Обоснуем формулу двумя способами:
        \begin{itemize}
            \item \emph{Через сочетания с повторениями (мультимножества):} 
            Любой неупорядоченный набор с возможными повторениями из $m$ 
            чисел, выбранных из множества $\{1, \dots, n\}$ (то есть 
            $m$-мультимножество), после сортировки по неубыванию 
            однозначно задаёт значения $f(1) \le f(2) \le \dots \le f(m)$. 
            Следовательно, число таких функций равно числу сочетаний с 
            повторениями из $n$ типов элементов по $m$:
            \begin{equation} \label{eq:map_weak_inc}
                N_{\text{нестр}} = \bar{C}_n^m = C_{n+m-1}^m = C_{n+m-1}^{n-1}. 
                \tag{6}
            \end{equation}
            \item \emph{Через сведение к строго возрастающим отображениям:} 
            Поставим каждому неубывающему отображению $f$ в соответствие 
            новое отображение $g(k) = f(k) + k - 1$ ($k = \overline{1, m}$). 
            Тогда из нестрогих неравенств $f(k) \le f(k + 1)$ следуют 
            строгие неравенства $g(k) < g(k + 1)$, причём 
            $1 \le g(1) < g(2) < \dots < g(m) \le n + m - 1$. По пункту 2 
            число таких отображений $g$ равно $C_{n+m-1}^m$.
        \end{itemize}
    \end{quote}
\end{proof}

\subsection{Задачи о подмножествах и диагоналях многоугольника}

\begin{problem}[\textbf{Пары непересекающихся подмножеств}]
    Дано конечное множество $U$ мощности $\abs{U} = n$. Найти число 
    упорядоченных пар подмножеств $(X, Y)$ множества $U$ ($X, Y \subseteq U$), 
    таких что $X \cap Y = \varnothing$.
\end{problem}

\begin{proof}[Решение (два способа)]~\\
    \begin{quote}
        \textbf{Способ 1 (разбиение по мощности множества $X$ и бином 
        Ньютона):}\\
        Пусть первое подмножество $X \subseteq U$ содержит ровно $k$ 
        элементов ($\abs{X} = k$, где $0 \le k \le n$). Выбрать такое 
        подмножество $X$ можно $C_n^k$ способами.

        Условие $X \cap Y = \varnothing$ означает, что второе подмножество 
        $Y$ не может содержать ни одного из $k$ элементов множества $X$, то 
        есть $Y$ является произвольным подмножеством дополнения 
        $U \setminus X$ мощности $\abs{U \setminus X} = n - k$. 
        Следовательно, при фиксированном $X$ подмножество $Y \subseteq U \setminus X$ 
        можно выбрать $2^{n-k}$ способами. Суммируя по всем возможным 
        значениям $k = 0, 1, \dots, n$ и сворачивая сумму по формуле 
        бинома Ньютона, получаем:
        \begin{equation} \label{eq:disjoint_subsets_binom}
            N = \sum_{k=0}^{n} C_n^k \cdot 2^{n-k} = 
            \sum_{k=0}^{n} C_n^k \cdot 1^k \cdot 2^{n-k} = (1 + 2)^n = 3^n. 
            \tag{7}
        \end{equation}

        \textbf{Способ 2 (комбинаторный выбор состояния для каждого 
        элемента):}\\
        Для каждого из $n$ элементов $u \in U$ при построении пары 
        непересекающихся подмножеств $(X, Y)$ имеется ровно $3$ 
        взаимоисключающие возможности:
        \begin{enumerate}
            \item $u \in X$ и $u \notin Y$ (элемент попадает только в $X$);
            \item $u \notin X$ и $u \in Y$ (элемент попадает только в $Y$);
            \item $u \notin X$ и $u \notin Y$ (элемент не попадает ни в $X$, 
            ни в $Y$).
        \end{enumerate}
        Четвёртый случай ($u \in X$ и $u \in Y$) запрещён условием 
        $X \cap Y = \varnothing$. Так как судьба каждого из $n$ элементов 
        выбирается независимо $3$ способами, по правилу умножения имеется 
        ровно $\underbrace{3 \cdot 3 \cdot \dots \cdot 3}_{n} = 3^n$ пар.
    \end{quote}
\end{proof}

\begin{problem}[\textbf{Непересекающиеся диагонали выпуклого $n$-угольника}]
    Дан выпуклый $n$-угольник ($n \ge 4$). Сколькими способами в нём можно 
    выбрать две непересекающиеся диагонали (не имеющие общих точек ни 
    внутри многоугольника, ни в его вершинах)? Порядок выбора диагоналей 
    неважен (пара $\{AB, CD\}$ считается один раз).
\end{problem}

\begin{proof}[Решение (Способ 1: метод вычитания из общего числа пар)]~\\
    \begin{quote}
        \textbf{1. Общее число пар диагоналей:}\\
        В выпуклом $n$-угольнике из каждой из $n$ вершин можно провести 
        $n - 3$ диагонали (во все вершины, кроме самой себя и двух соседних). 
        Так как каждая диагональ соединяет две вершины, общее число 
        диагоналей равно:
        \[
            D = \frac{n(n - 3)}{2}.
        \]
        Общее количество неупорядоченных пар различных диагоналей равно 
        числу сочетаний $C_D^2 = C_{\frac{n(n-3)}{2}}^2$.

        \textbf{2. Исключение пересекающихся пар диагоналей:}\\
        Две различные диагонали выпуклого $n$-угольника имеют общую точку 
        в одном из двух непересекающихся случаев:
        \begin{itemize}
            \item \emph{Диагонали пересекаются строго внутри $n$-угольника:} 
            Любые две такие диагонали однозначно определяют $4$ вершины 
            $n$-угольника (свои концы), и наоборот: любые $4$ вершины 
            выпуклого $n$-угольника образуют выпуклый четырёхугольник, 
            диагонали которого пересекаются строго внутри $n$-угольника. 
            Следовательно, число пар внутренне пересекающихся диагоналей 
            в точности равно $C_n^4$.
            \item \emph{Диагонали имеют общую вершину (пересекаются на 
            границе):} 
            В каждой из $n$ вершин многоугольника сходятся ровно $n - 3$ 
            диагонали. Выбрать две диагонали, выходящие из одной 
            фиксированной вершины, можно $C_{n-3}^2$ способами. Поскольку 
            две различные диагонали не могут иметь более одной общей 
            вершины, суммируя по всем $n$ вершинам, получаем $n \cdot C_{n-3}^2$ 
            пар.
        \end{itemize}
        Вычитая из общего числа пар оба запрещённых случая, получаем:
        \begin{equation} \label{eq:diag_sub_method}
            N = C_{\frac{n(n-3)}{2}}^2 - C_n^4 - n \cdot C_{n-3}^2. \tag{8}
        \end{equation}
    \end{quote}
\end{proof}

\begin{proof}[Решение (Способ 2: прямое комбинаторное решение через четвёрки вершин)]~\\
    \begin{quote}
        Ответим на вопрос из конспекта о существовании другого (прямого) 
        решения. Поскольку две искомые диагонали не имеют общих вершин, их 
        концы образуют четвёрку попарно различных вершин $A, B, C, D$ 
        $n$-угольника (занумеруем их в порядке обхода по периметру).

        \begin{figure}[H]
            \centering
            \begin{tikzpicture}[scale=1.1]
                % Схема четырёх вершин A, B, C, D
                \coordinate (A) at (0, 0);
                \coordinate (B) at (0.8, 1.3);
                \coordinate (D) at (1.5, -0.2);
                \coordinate (C) at (2.3, 1.1);

                \draw[blue!80!black, ultra thick] (A) -- (B);
                \draw[blue!80!black, ultra thick] (D) -- (C);
                \draw[gray, dashed] (B) -- (C);
                \draw[gray, dashed] (A) -- (D);

                \fill (A) circle (2pt) node[below left] {$A$};
                \fill (B) circle (2pt) node[above left] {$B$};
                \fill (C) circle (2pt) node[above right] {$C$};
                \fill (D) circle (2pt) node[below right] {$D$};

                \node[right, align=left] at (3.0, 0.55) {
                    Пара $\{AB, CD\}$ считается $1$ раз\\
                    (противоположные стороны $ABCD$,\\
                    не являющиеся сторонами $n$-угольника)
                };
            \end{tikzpicture}
            \caption{Две непересекающиеся диагонали образуют пару 
            противоположных сторон вписанного четырёхугольника $ABCD$.}
        \end{figure}

        На любых $4$ вершинах $A, B, C, D$ можно построить ровно $2$ пары 
        непересекающихся отрезков: $\{AB, CD\}$ и $\{BC, AD\}$ (третья пара 
        $\{AC, BD\}$ пересекается внутри). Всего таких пар непересекающихся 
        хорд в $n$-угольнике $2 C_n^4$. Исключим из них те пары, в которых 
        хотя бы один отрезок оказался \emph{стороной} $n$-угольника:
        \begin{itemize}
            \item Пары из \emph{двух несмежных сторон} $n$-угольника: из 
            всех $C_n^2$ пар сторон вычитаем $n$ пар смежных сторон, получая 
            $C_n^2 - n = \frac{n(n - 3)}{2}$ пар.
            \item Пары вида \emph{«одна сторона + одна непересекающаяся с 
            ней диагональ»}: сторону можно выбрать $n$ способами; после 
            удаления $2$ её концов остаётся $n - 2$ вершины, образующие 
            $C_{n-2}^2$ отрезков, среди которых ровно $n - 3$ являются 
            сторонами $n$-угольника, а остальные $C_{n-2}^2 - (n - 3)$ --- 
            диагоналями. Всего таких пар $n\bigl(C_{n-2}^2 - (n - 3)\bigr)$.
        \end{itemize}
        Вычитая эти пары из $2 C_n^4$, получаем вторую формулу:
        \begin{equation} \label{eq:diag_direct_method}
            N = 2 C_n^4 - n\bigl(C_{n-2}^2 - n + 3\bigr) - \frac{n(n - 3)}{2}. 
            \tag{9}
        \end{equation}
        Раскрывая факториалы в любой из формул \eqref{eq:diag_sub_method} 
        или \eqref{eq:diag_direct_method}, приходим к компактному 
        замкнутому выражению (при $n \ge 6$, а для $n = 4, 5$ формула 
        корректно даёт $0$):
        \begin{equation} \label{eq:diag_closed_form}
            N = \frac{n(n - 3)(n - 4)(n - 5)}{12} = \frac{1}{2} n C_{n-3}^3. 
            \tag{10}
        \end{equation}
    \end{quote}
\end{proof}

\subsection{Задача о беспорядках ($n$ шляп и $n$ гостей)}

\begin{definition}
    \textdef{Беспорядком} (или \textdef{перестановкой без неподвижных точек}) 
    называется такая перестановка $\pi: \{1, \dots, n\} \to \{1, \dots, n\}$, 
    при которой ни один элемент не переходит сам в себя:
    \[
        \forall i \in \{1, 2, \dots, n\} : \pi(i) \neq i.
    \]
    Количество всех беспорядков на $n$ элементах обозначается через $D_n$ 
    (или $!n$ --- \textdef{субфакториал} числа $n$).
\end{definition}

\begin{problem}[\textbf{Задача о шляпах}]
    В гардеробе $n$ гостей сдали свои шляпы. Уходя домой, они в темноте 
    разбирают шляпы случайным образом (каждый берёт ровно по одной шляпе). 
    Сколькими способами гости могут разобрать шляпы так, чтобы \textbf{ни 
    один} гость не получил свою собственную шляпу? Решить задачу двумя 
    способами:
    \begin{enumerate}
        \item методом пересечений (через принцип включений и исключений);
        \item рекуррентным способом.
    \end{enumerate}
\end{problem}

\begin{proof}[Решение (Способ 1: через пересечения и формулу включений-исключений)]~\\
    \begin{quote}
        Обозначим через $S_n$ множество всех $n!$ перестановок $n$ шляп 
        между $n$ гостями. Для каждого $i \in \{1, 2, \dots, n\}$ введём 
        множество $A_i \subseteq S_n$, состоящее из всех перестановок, в 
        которых $i$-й гость получил \emph{свою} шляпу (то есть $i$ является 
        неподвижной точкой):
        \[
            A_i = \{\pi \in S_n \mid \pi(i) = i\}.
        \]
        Тогда объединение $\bigcup_{i=1}^n A_i$ содержит все перестановки, 
        в которых \emph{хотя бы один} гость получил свою шляпу, а число 
        беспорядков $D_n$ равно мощности дополнения этого объединения:
        \[
            D_n = \abs{S_n \setminus \bigcup_{i=1}^{n} A_i} = 
            n! - \abs{\bigcup_{i=1}^{n} A_i}.
        \]
        Найдём мощность пересечения любых $k$ множеств 
        $A_{i_1} \cap A_{i_2} \cap \dots \cap A_{i_k}$ (где 
        $1 \le i_1 < i_2 < \dots < i_k \le n$). Перестановка принадлежит 
        этому пересечению тогда и только тогда, когда гости с номерами 
        $i_1, \dots, i_k$ получили свои шляпы ($\pi(i_s) = i_s$), а 
        оставшиеся $n - k$ шляп распределены между остальными $n - k$ 
        гостями произвольным образом, что даёт $(n - k)!$ вариантов:
        \[
            \abs{A_{i_1} \cap A_{i_2} \cap \dots \cap A_{i_k}} = (n - k)!.
        \]
        Поскольку выбрать $k$ индексов $1 \le i_1 < \dots < i_k \le n$ 
        можно $C_n^k$ способами, сумма мощностей всех $k$-кратных 
        пересечений равна:
        \[
            \sum_{1 \le i_1 < \dots < i_k \le n} 
            \abs{A_{i_1} \cap \dots \cap A_{i_k}} = 
            C_n^k \cdot (n - k)! = \frac{n!}{k! \, \cancel{(n - k)!}} \cdot 
            \cancel{(n - k)!} = \frac{n!}{k!}.
        \]
        Применяя принцип включений и исключений, находим число 
        беспорядков:
        \begin{equation} \label{eq:derange_pie}
            D_n = n! - \sum_{k=1}^{n} (-1)^{k-1} \frac{n!}{k!} = 
            n! \sum_{k=0}^{n} \frac{(-1)^k}{k!} = 
            n!\left(1 - \frac{1}{1!} + \frac{1}{2!} - \frac{1}{3!} + \dots + 
            \frac{(-1)^n}{n!}\right). \tag{11}
        \end{equation}
        В частности, при $n \to \infty$ сумма в скобках сходится к 
        разложению Тейлора экспоненты $e^{-1} = \sum_{k=0}^\infty \frac{(-1)^k}{k!}$, 
        поэтому $D_n \sim \frac{n!}{e}$, а вероятность беспорядка 
        стремится к $\frac{1}{e} \approx 0,3679$.
    \end{quote}
\end{proof}

\begin{proof}[Решение (Способ 2: вывод рекуррентного соотношения)]~\\
    \begin{quote}
        Проследим за судьбой $1$-го гостя и $1$-й шляпы при $n \ge 2$ 
        (причём $D_1 = 0$, $D_2 = 1$). Так как $1$-й гость не может взять 
        свою шляпу №$1$, он берёт некоторую чужую шляпу с номером 
        $k \in \{2, 3, \dots, n\}$ ($\pi(1) = k$). Выбрать такой номер $k$ 
        можно ровно \textbf{$(n - 1)$ способом}.

        После того как $1$-й гость забрал шляпу №$k$, дальнейшее 
        распределение распадается на два непересекающихся случая в 
        зависимости от того, какую шляпу возьмёт $k$-й гость:
        \begin{enumerate}
            \item \textbf{Случай А: $k$-й гость берёт шляпу №$1$ 
            ($\pi(k) = 1$).} 
            В этом случае гости $1$ и $k$ просто обменялись шляпами друг с 
            другом (образовали транспозицию --- цикл длины $2$). Тогда 
            оставшиеся $n - 2$ гостей должны разобрать свои $n - 2$ шляпы 
            так, чтобы никто не получил свою, что по определению можно 
            сделать ровно $D_{n-2}$ способами.
            \item \textbf{Случай Б: $k$-й гость НЕ берёт шляпу №$1$ 
            ($\pi(k) \neq 1$).} 
            Тогда для оставшихся $n - 1$ гостей (с номерами 
            $\{2, 3, \dots, n\}$) остаются $n - 1$ шляп (с номерами 
            $\{1, 2, \dots, n\} \setminus \{k\}$). При этом каждому 
            $j$-му гостю ($j \neq k$) запрещено брать свою шляпу №$j$, а 
            $k$-му гостю запрещено брать шляпу №$1$ (она играет для него 
            роль «запретной своей» шляпы!). Мы получили задачу о 
            беспорядках для $n - 1$ гостей и $n - 1$ шляп, которая имеет 
            ровно $D_{n-1}$ решений.
        \end{enumerate}
        Объединяя оба случая по правилу сложения и умножая на $(n - 1)$ 
        вариантов выбора гостя $k$, получаем линейное рекуррентное 
        соотношение второго порядка:
        \begin{equation} \label{eq:derange_rec_1}
            D_n = (n - 1)(D_{n-1} + D_{n-2}), \qquad D_0 = 1, \ D_1 = 0. \tag{12}
        \end{equation}
        Перепишем соотношение \eqref{eq:derange_rec_1} в виде разности:
        \[
            D_n - n D_{n-1} = -(D_{n-1} - (n - 1)D_{n-2}).
        \]
        Отсюда по индукции следует, что выражение $D_n - n D_{n-1}$ на 
        каждом шаге меняет знак и равно $(-1)^{n-1}(D_1 - 1 \cdot D_0) = (-1)^n$, 
        что даёт рекуррентное соотношение первого порядка:
        \begin{equation} \label{eq:derange_rec_2}
            D_n = n D_{n-1} + (-1)^n, \qquad n \ge 1. \tag{13}
        \end{equation}
        Последовательное вычисление по формуле \eqref{eq:derange_rec_1} 
        или \eqref{eq:derange_rec_2} даёт первые значения субфакториала: 
        $D_1 = 0$, $D_2 = 1$, $D_3 = 2$, $D_4 = 9$, $D_5 = 44$, $D_6 = 265$.
    \end{quote}
\end{proof}